\vssub
\subsection{~引言} \label{sec:intro}
\vssub

在有限水深且平均流不为零的水体中，波浪或谱波分量通常用若干相位参数和
振幅参数描述。相位参数包括波数向量 {\bk}、波数 $k$、方向
$\theta$ 以及若干频率。若需要考虑平均流对波浪的影响，则要区分相对频率
或本征（弧度）频率 $\sigma$ $(= 2 \pi f_r)$ 与绝对（弧度）频率
$\omega$ $(= 2 \pi f_a)$：前者是在随平均流运动的参考系中观测到的，
后者是在固定参考系中观测到的。按定义，方向 $\theta$ 垂直于波峰
（或谱分量的波峰），并且与 {\bk} 的方向一致。这里给出的方程遵循几何光学
近似；当水深和流场变化尺度远大于单个波浪尺度时，该近似在极限意义下是
精确的\footnote{即使波高在半个波长内变化达 5 倍，几何光学近似仍可给出
合理结果，这一点已在海底峡谷算例中得到说明 \citep{art:Mea07}。}。该近似
忽略的绕射、散射和干涉效应，可以在事后作为源项加入波作用量方程。在这种
缓变流和缓变水深近似下，可在局地采用准均匀（线性）波理论，从而得到下面
的色散关系和多普勒型方程，用以关联各相位参数：

%-----------------------------------%
% dispersion and Doppler equations  %
%-----------------------------------%
% eq:disp
% eq:doppler

\begin{equation}
\sigma ^2 = g k \tanh kd \: ,
\label{eq:disp}
\end{equation}
\begin{equation}
\omega = \sigma + {\bk} \cdot {\bf U} \: ,
\label{eq:doppler}
\end{equation}

\noindent 
其中 $d$ 为平均水深，{\bf U} 为流速（在单个波浪尺度上对水深和时间取平均）。
水深和流场缓变的假设意味着海底地形具有大尺度变化，因此通常可以忽略波浪
绕射。由波浪或波分量的相位函数给出 {\bk} 和 $\omega$ 的通常定义，意味着
波峰数守恒 \citep[参见，例如][]{bk:Phi77,bk:Mei83}：

%-----------------------------%
% conservation of wave crests %
%-----------------------------%
% eq:wnconv

\begin{equation}
\frac{\partial {\bk}}{\partial t} + \nabla \omega = 0 \: .
\label{eq:wnconv}
\end{equation}

%\noindent

由式~(\ref{eq:disp}) 至式~(\ref{eq:wnconv}) 可以计算相位参数的变化率
\citep[例如，相关方程此处不再列出][]{rep:Chr82,bk:Mei83,tol:JPO90}。

%For monochromatic waves, the amplitude is described as the amplitude, the wave
%height, or the wave energy. 
对于不规则风浪，海面（随机）方差使用表面高程方差密度谱 $F$ 描述。在波浪
模拟界，这通常称为“能量谱”。该谱 $F$ 是所有独立相位参数的函数，即
$F({\bk},\sigma,\omega)$；此外，它还在大于单个波浪尺度的空间和时间尺度上
变化，例如 $F({\bk},\sigma,\omega;\bx,t)$。然而，对于频率高于主导风海频率
3 倍的短波 \citep{Leckler&al.2015}，能量与线性色散关系非常接近，因此
式~(\ref{eq:disp}) 和式~(\ref{eq:doppler}) 将 {\bk}、$\sigma$ 与
$\omega$ 关联起来。因此只有两个独立相位参数，局地瞬时谱成为二维谱。
束缚波非线性贡献的影响可在 \ws\ 中按 \cite{art:Jan09} 的方法作为后处理加入。

在 \ws\ 中，基本谱为波数--方向谱 $F(k,\theta)$。之所以选择该谱，是因为
在可变水深下，就波浪增长和衰减的物理过程而言，它具有不变性特征。然而
\ws\ 的输出采用更传统的频率--方向谱 $F(f_r,\theta)$。可通过直接的
Jacobian 变换，由 $F(k,\theta)$ 计算不同形式的谱：

%--------------------------------------------%
% Jacobian transformation and group velocity %
%--------------------------------------------%
% eq:jac_fr
% eq:jac_fa
% eq:cg

\begin{equation}
F(f_r,\theta) = \frac{\partial k}{\partial f_r} F(k,\theta) =
\frac{2\pi}{c_g} F(k,\theta) \: ,
\label{eq:jac_fr}
\end{equation}
\begin{equation}
F(f_a,\theta) = \frac{\partial k}{\partial f_a} F(k,\theta) =
\frac{2\pi}{c_g}
\left ( 1 + \frac{{\bk}\cdot{\bf U}}{k c_g}\right )^{-1}
F(k,\theta) \: ,
\label{eq:jac_fa}
\end{equation}
\begin{equation}
c_g = \frac{\partial \sigma}{\partial k} = n \frac{\sigma}{k}
\; , \;
n = \frac{1}{2} + \frac{kd}{\sinh 2kd} \; ,
\label{eq:cg}
\end{equation}

\noindent 
其中 $c_g$ 为群速。由这些谱中的任意一种，都可以通过对方向积分生成一维谱；
而按定义，对整个谱积分得到总方差 $E$（在波浪模拟界通常称为波能）。

在无流情况下，波包的方差（能量）是守恒量。在有流情况下，由于流对波浪平均
动量传递做功，某一谱分量的能量或方差不再守恒
\citep{art:LHS61,art:LHS62}。然而，从一般意义上说，波作用量
$A \equiv E/\sigma$ 是守恒的 \citep[例如][]{art:Whi65,art:BG68}。
因此，波作用量密度谱 $N(k,\theta) \equiv F(k,\theta)/\sigma$ 是本模型中
优先采用的谱。于是波浪传播可描述为

%--------------------------------------%
% Most general action balance equation %
%--------------------------------------%
% eq:balance0

\begin{equation}
\frac{D N}{D t} = \frac{S}{\sigma} \: ,
\label{eq:balance0}
\end{equation}

\noindent 
其中 $D/Dt$ 表示全导数（随波分量运动），$S$ 表示谱 $F$ 的源项和汇项的
净效应。由于式~(\ref{eq:balance0}) 左端通常描述不含散射的线性传播，
非线性波浪传播效应（即波--波相互作用）和部分波浪反射体现在 $S$ 中。
传播项和源项将在以下各节分别讨论。
